\documentclass[12pt,a4paper]{article}

\usepackage[margin=2.5cm]{geometry}
\usepackage{hyperref}
\usepackage{listings}
\usepackage{xcolor}
\usepackage{booktabs}
\usepackage{enumitem}
\usepackage{parskip}
\usepackage{graphicx}
\usepackage{titlesec}
\usepackage{mdframed}

\hypersetup{colorlinks=true, linkcolor=blue, urlcolor=blue}

\definecolor{cmdbg}{gray}{0.93}
\lstdefinestyle{cmd}{
  backgroundcolor=\color{cmdbg},
  basicstyle=\ttfamily\small,
  breaklines=true,
  frame=single,
  frameround=tttt,
  framesep=4pt,
  xleftmargin=6pt,
  xrightmargin=6pt,
}
\lstset{style=cmd}

\newmdenv[
  backgroundcolor=yellow!15,
  linecolor=orange!70!black,
  linewidth=1pt,
  frametitle={\small\textbf{Note}},
  frametitlefont=\bfseries\small,
  innertopmargin=6pt,
  innerbottommargin=6pt,
]{note}

\title{\textbf{Recipe Maker — Windows Setup Guide}\\[0.4em]
  \large Step-by-step instructions to run the full application on Windows}
\author{Software Engineering — Academic Year 2025/2026}
\date{4 May 2026}

\begin{document}
\maketitle
\tableofcontents
\newpage

% ─────────────────────────────────────────────────────────────
\section{Prerequisites}

The following software must be installed on your Windows machine before
attempting to run the project.  All tools listed below are free.

\subsection{Required Software}

\begin{center}
\begin{tabular}{llll}
\toprule
\textbf{Tool} & \textbf{Min.\ version} & \textbf{Purpose} & \textbf{Download URL} \\
\midrule
JDK (Java) & 21 & Backend runtime / compiler &
  \href{https://adoptium.net}{adoptium.net} \\
Node.js    & 20 LTS & Frontend runtime / npm &
  \href{https://nodejs.org}{nodejs.org} \\
Git        & 2.x & Clone the repository &
  \href{https://git-scm.com}{git-scm.com} \\
\bottomrule
\end{tabular}
\end{center}

\begin{note}
Maven is \textbf{not} required to be installed globally.  The project ships with
the Maven Wrapper (\texttt{mvnw.cmd}) which downloads the correct Maven version
automatically on first use.
\end{note}

\subsection{Recommended (Optional)}

\begin{itemize}[noitemsep]
  \item \textbf{Visual Studio Code} (\href{https://code.visualstudio.com}{code.visualstudio.com})
        with the \emph{Extension Pack for Java} and \emph{ES7+ React/Redux/React-Native snippets}.
  \item \textbf{IntelliJ IDEA Community} (\href{https://www.jetbrains.com/idea}{jetbrains.com/idea})
        for a richer Java development experience.
  \item \textbf{Windows Terminal} for a modern command-line experience.
\end{itemize}

% ─────────────────────────────────────────────────────────────
\section{Installing the Prerequisites}

\subsection{Java 21 (JDK)}
\begin{enumerate}[noitemsep]
  \item Open \url{https://adoptium.net} in your browser.
  \item Select \textbf{Temurin 21 (LTS)}, OS: \emph{Windows}, Architecture: \emph{x64}.
  \item Download the \texttt{.msi} installer and run it.
  \item On the installer's \emph{Custom Setup} screen, make sure
        \textbf{Set JAVA\_HOME} and \textbf{Add to PATH} are both enabled (they
        are by default).
  \item After installation, open a new \textbf{Command Prompt} or
        \textbf{PowerShell} window and verify:
\begin{lstlisting}
java -version
\end{lstlisting}
        Expected output (version number may differ in the minor digits):
\begin{lstlisting}
openjdk version "21.x.x" ...
\end{lstlisting}
\end{enumerate}

\subsection{Node.js 20 LTS}
\begin{enumerate}[noitemsep]
  \item Open \url{https://nodejs.org} and download the \textbf{LTS} installer
        for Windows (\texttt{.msi}).
  \item Run the installer and accept all defaults.  Ensure \emph{Add to PATH}
        is checked.
  \item Verify the installation:
\begin{lstlisting}
node --version
npm --version
\end{lstlisting}
        Both commands should return a version number without errors.
\end{enumerate}

\subsection{Git}
\begin{enumerate}[noitemsep]
  \item Download the installer from \url{https://git-scm.com/download/win}.
  \item Run the installer.  The defaults are suitable; click \emph{Next}
        through all screens.
  \item Verify:
\begin{lstlisting}
git --version
\end{lstlisting}
\end{enumerate}

% ─────────────────────────────────────────────────────────────
\section{Obtaining the Source Code}

Open a \textbf{Command Prompt} or \textbf{PowerShell} window and navigate to
the directory where you want to place the project, then clone the repository:

\begin{lstlisting}
cd C:\Users\YourName\Projects
git clone <repository-url> recipe_maker
cd recipe_maker
\end{lstlisting}

Replace \texttt{<repository-url>} with the actual Git remote URL (HTTPS or
SSH) provided by your team or institution.

If you already have the project as a ZIP archive, extract it and
\texttt{cd} into the extracted folder instead.

% ─────────────────────────────────────────────────────────────
\section{Project Structure}

After cloning, the top-level layout is:

\begin{lstlisting}
recipe_maker/
  backend/          <- Spring Boot application (Java 21, Maven)
    mvnw.cmd        <- Maven Wrapper for Windows
    pom.xml         <- Maven build descriptor
    src/
  frontend/         <- React 18 + TypeScript + Vite application
    package.json
    vite.config.ts
    src/
  Dockerfile
  PROJECT_REPORT.tex
  SETUP_WINDOWS.tex <- this document
\end{lstlisting}

The backend exposes its REST API on \textbf{port 8081}.  The frontend
development server runs on \textbf{port 5173} and proxies all
\texttt{/api/**} requests to the backend automatically.

% ─────────────────────────────────────────────────────────────
\section{Running the Backend}

Open a \textbf{Command Prompt} or \textbf{PowerShell} window.

\subsection{Navigate to the backend directory}
\begin{lstlisting}
cd C:\Users\YourName\Projects\recipe_maker\backend
\end{lstlisting}

\subsection{Start the Spring Boot application}
\begin{lstlisting}
mvnw.cmd spring-boot:run
\end{lstlisting}

\begin{note}
On first run, Maven will download all dependencies from the internet
(\textasciitilde{}100--200 MB).  Subsequent starts are much faster because
the artifacts are cached in \texttt{C:\textbackslash Users\textbackslash
YourName\textbackslash .m2\textbackslash repository}.
\end{note}

The application is ready when you see a log line similar to:
\begin{lstlisting}
Started RecipeMakerApplication in X.XXX seconds
\end{lstlisting}

\subsection{Verify the backend}
Open a browser and go to:
\begin{itemize}[noitemsep]
  \item \url{http://localhost:8081/swagger-ui.html} — interactive API
        documentation (Swagger UI).
  \item \url{http://localhost:8081/h2-console} — H2 in-memory database
        console.  JDBC URL: \texttt{jdbc:h2:mem:recipemaker}, username:
        \texttt{sa}, password: \emph{(empty)}.
\end{itemize}

\subsection{Running the tests (optional)}
\begin{lstlisting}
mvnw.cmd test
\end{lstlisting}
Test reports are written to
\texttt{backend\textbackslash target\textbackslash surefire-reports\textbackslash}.

% ─────────────────────────────────────────────────────────────
\section{Running the Frontend}

Open a \textbf{second} Command Prompt or PowerShell window (keep the backend
running in the first one).

\subsection{Navigate to the frontend directory}
\begin{lstlisting}
cd C:\Users\YourName\Projects\recipe_maker\frontend
\end{lstlisting}

\subsection{Install dependencies (first time only)}
\begin{lstlisting}
npm install
\end{lstlisting}

This downloads all Node.js packages listed in \texttt{package.json} into a
local \texttt{node\_modules} folder (\textasciitilde{}150--250 MB).

\subsection{Start the development server}
\begin{lstlisting}
npm run dev
\end{lstlisting}

Expected output:
\begin{lstlisting}
  VITE v5.x.x  ready in XXX ms

  Local:   http://localhost:5173/
  Network: use --host to expose
\end{lstlisting}

Open \url{http://localhost:5173} in your browser to access the application.

\subsection{Building for production (optional)}
\begin{lstlisting}
npm run build
\end{lstlisting}

Static files are written to \texttt{frontend\textbackslash dist\textbackslash}.

% ─────────────────────────────────────────────────────────────
\section{Default Accounts}

The backend seeds several demo accounts on startup.  Use any of the
following credentials to log in:

\begin{center}
\begin{tabular}{lll}
\toprule
\textbf{Username} & \textbf{Password} & \textbf{Role} \\
\midrule
\texttt{admin}    & \texttt{admin}    & ADMIN \\
\texttt{doctor1}  & \texttt{password} & DOCTOR \\
\texttt{dietician1} & \texttt{password} & DIETICIAN \\
\texttt{patient1} & \texttt{password} & PATIENT \\
\texttt{kitchen1} & \texttt{password} & KITCHEN \\
\bottomrule
\end{tabular}
\end{center}

\begin{note}
Because the database is in-memory (\texttt{H2 mem}), all data is lost when
the backend process stops.  Seed data is re-inserted automatically on the
next startup.
\end{note}

% ─────────────────────────────────────────────────────────────
\section{Startup Sequence Summary}

Always start the services in this order:

\begin{enumerate}
  \item \textbf{Backend first} — Spring Boot must be running before the
        frontend can make API calls.
  \item \textbf{Frontend second} — Vite dev server proxies API calls to the
        backend.
\end{enumerate}

\begin{center}
\begin{tabular}{lll}
\toprule
\textbf{Terminal} & \textbf{Directory} & \textbf{Command} \\
\midrule
1 (backend)  & \texttt{recipe\_maker\textbackslash backend}  & \texttt{mvnw.cmd spring-boot:run} \\
2 (frontend) & \texttt{recipe\_maker\textbackslash frontend} & \texttt{npm run dev} \\
\bottomrule
\end{tabular}
\end{center}

% ─────────────────────────────────────────────────────────────
\section{Ports and URLs Reference}

\begin{center}
\begin{tabular}{lll}
\toprule
\textbf{Service} & \textbf{URL} & \textbf{Notes} \\
\midrule
Frontend (dev)   & \url{http://localhost:5173}              & Main app UI \\
Backend API      & \url{http://localhost:8081/api}          & REST endpoints \\
Swagger UI       & \url{http://localhost:8081/swagger-ui.html} & API explorer \\
OpenAPI JSON     & \url{http://localhost:8081/v3/api-docs}  & Raw spec \\
H2 Console       & \url{http://localhost:8081/h2-console}   & DB browser \\
\bottomrule
\end{tabular}
\end{center}

% ─────────────────────────────────────────────────────────────
\section{Troubleshooting}

\subsection{\texttt{mvnw.cmd} is not recognised}
Ensure you are inside the \texttt{backend} directory.  The wrapper script
must be executed from there.  Also confirm Java 21 is on the PATH:
\begin{lstlisting}
java -version
\end{lstlisting}

\subsection{Port 8081 already in use}
Another process is occupying the port.  Find and stop it:
\begin{lstlisting}
netstat -ano | findstr :8081
taskkill /PID <PID> /F
\end{lstlisting}
Alternatively, change the backend port in
\texttt{backend\textbackslash src\textbackslash main\textbackslash resources\textbackslash
application.properties}:
\begin{lstlisting}
server.port=9090
\end{lstlisting}
and update \texttt{vite.config.ts} accordingly:
\begin{lstlisting}[language=Java]
proxy: {
  '/api': { target: 'http://localhost:9090', changeOrigin: true }
}
\end{lstlisting}

\subsection{Port 5173 already in use}
Vite will automatically try the next available port and print the actual URL.
Alternatively specify a port explicitly:
\begin{lstlisting}
npm run dev -- --port 5174
\end{lstlisting}

\subsection{\texttt{npm install} fails with network errors}
If you are behind a corporate proxy, configure npm to use it:
\begin{lstlisting}
npm config set proxy http://proxy.company.com:PORT
npm config set https-proxy http://proxy.company.com:PORT
\end{lstlisting}
For Maven, create or edit
\texttt{C:\textbackslash Users\textbackslash YourName\textbackslash
.m2\textbackslash settings.xml} and add a \texttt{<proxies>} section
following the official Maven proxy guide.

\subsection{CORS or 401 errors in the browser}
Ensure the backend is running before loading the frontend, and that you are
accessing the app via \texttt{http://localhost:5173} (not via a file path).
The Vite proxy only works for requests made from the dev server origin.

\subsection{Execution Policy error running \texttt{mvnw.cmd}}
If PowerShell blocks the script, use \textbf{Command Prompt} (\texttt{cmd.exe})
instead, or run:
\begin{lstlisting}
Set-ExecutionPolicy -Scope CurrentUser -ExecutionPolicy RemoteSigned
\end{lstlisting}

% ─────────────────────────────────────────────────────────────
\section{Firewall and Antivirus}
Windows Defender may prompt to allow network access for Java and Node.js.
Click \emph{Allow access} for both when prompted.  No inbound firewall rules
are needed for local development.

% ─────────────────────────────────────────────────────────────
\section{Quick-Start Checklist}

\begin{itemize}
  \item[$\square$] JDK 21 installed and \texttt{java -version} returns 21.x.x
  \item[$\square$] Node.js 20 LTS installed and \texttt{node --version} returns 20.x.x
  \item[$\square$] Git installed and repository cloned
  \item[$\square$] Terminal 1: \texttt{cd backend} then \texttt{mvnw.cmd spring-boot:run}
  \item[$\square$] Backend log shows ``Started RecipeMakerApplication''
  \item[$\square$] Terminal 2: \texttt{cd frontend} then \texttt{npm install} (first time)
  \item[$\square$] Terminal 2: \texttt{npm run dev}
  \item[$\square$] Browser open at \url{http://localhost:5173}
  \item[$\square$] Log in with \texttt{admin} / \texttt{admin}
\end{itemize}

\end{document}
